---
## Front matter
title: "Отчёт по лабораторной работе №8"
subtitle: "Computer Skills for Scientific Writing"
author: "Сунь Маосин"

## Generic otions
lang: ru-RU
toc-title: "Содержание"

## Pdf output format
toc: true
toc-depth: 2
lof: true
lot: true
fontsize: 12pt
linestretch: 1.5
papersize: a4
documentclass: scrreprt
## I18n polyglossia
polyglossia-lang:
  name: russian
  options:
    - spelling=modern
    - babelshorthands=true
polyglossia-otherlangs:
  name: english
## I18n babel
babel-lang: russian
babel-otherlangs: english
## Fonts
mainfont: Times New Roman
romanfont: Times New Roman
sansfont: Arial
monofont: Courier New
mathfont: Times New Roman
mainfontoptions: Ligatures=Common,Ligatures=TeX,Scale=0.94
romanfontoptions: Ligatures=Common,Ligatures=TeX,Scale=0.94
sansfontoptions: Ligatures=Common,Ligatures=TeX,Scale=MatchLowercase,Scale=0.94
monofontoptions: Scale=MatchLowercase,Scale=0.94,FakeStretch=0.9
mathfontoptions:
## Biblatex
biblatex: true
biblio-style: "gost-numeric"
biblatexoptions:
  - parentracker=true
  - backend=biber
  - hyperref=auto
  - language=auto
  - autolang=other*
  - citestyle=gost-numeric
## Pandoc-crossref LaTeX customization
figureTitle: "Рис."
tableTitle: "Таблица"
listingTitle: "Листинг"
lofTitle: "Список иллюстраций"
lotTitle: "Список таблиц"
lolTitle: "Листинги"
## Misc options
indent: true
header-includes:
  - \usepackage{indentfirst}
  - \usepackage{float}
  - \floatplacement{figure}{H}
---

# Цель работы

Изучение возможностей пакета TikZ для создания векторной графики в LaTeX.

# Ход выполнения

## Компиляция исходного файла

### Компиляция

![Компиляция tikz_drawing.tex](image/01.png)

На первом этапе был открыт исходный файл `tikz_drawing.tex` и выполнена его компиляция с помощью команды `pdflatex`.

## Основные геометрические фигуры

### Код

![Код геометрических фигур](image/02.png)

### Результат

![Геометрические фигуры](image/03.png)

В ходе упражнения были созданы различные типы линий (прямые, со стрелками), прямоугольники (с заливкой и без), окружности и эллипсы.

## Работа с цветом и стилями

### Код

![Код цветных линий](image/04.png)

### Результат

![Цветные линии](image/05.png)

Протестированы различные цвета (красный, синий, зелёный, оранжевый, фиолетовый) и стили линий (сплошная, пунктирная, штрихпунктирная).

## Простые узлы

### Код

![Код простых узлов](image/06.png)

### Результат

![Простые узлы](image/07.png)

Созданы узлы различных типов: без рамки, с рамкой, круглые, прямоугольные с заливкой.

## Соединение узлов

### Код

![Код соединения узлов](image/08.png)

### Результат

![Соединение узлов](image/09.png)

Реализовано соединение узлов прямыми и изогнутыми линиями, а также построена цепочка узлов с ответвлением.

## Графики функций

### Код

![Код графиков функций](image/10.png)

### Результат

![Графики функций](image/11.png)

Построены графики квадратичной функции, а также тригонометрических функций синуса и косинуса.

## Использование циклов

### Код

![Код с циклами](image/12.png)

### Результат

![Результат использования циклов](image/13.png)

С помощью циклов `foreach` созданы повторяющиеся элементы (окружности) и координатная сетка с подписями.

## Сложные примеры

В работе также были созданы более сложные структуры: простой граф с тремя узлами и схема алгоритма с использованием различных стилей узлов.

# Вывод

В ходе выполнения лабораторной работы были изучены основные возможности пакета TikZ:

- создание геометрических фигур (линии, прямоугольники, окружности, эллипсы);
- работа с цветами и стилями линий;
- создание и соединение узлов;
- построение графиков функций;
- использование циклов для повторяющихся элементов;
- создание сложных схем и графов.

Все файлы успешно скомпилированы, полученный PDF-документ полностью соответствует ожидаемым результатам.